% !TeX root = ../papers/nc-p3.tex
\section{三文字の可換部分と次数七の点グラフ}
\label{sec:septic-point-curves}
\begin{proposition}\label{prop:commuting-three-point-graph}
有限双対が標準Frobenius型で、各辺の左右の射影ファイバーが空または一点とする。
三次元空間 $H$ と各位置での単射 $H\hookrightarrow V_i$ があり、
全辺で $\bigwedge^2 H\subset R_i$ なら、元代数は標準Koszul左右AS正則である。
可換三変数多項式環の反復 $P$ はこれらの文字で埋め込まれ、元代数は左右 $P$ 自由である。
三文字の直交点を次の位置の三文字の直交点へ運ぶという条件は不要である。
\end{proposition}
\begin{proof}
位置 $i$ の三文字を消す点 $p$ で関係を評価すると、三つの可換関係が零になるから階数は高々3であり、
ファイバー条件でちょうど3である。従って $p$ は点グラフ上にあり、前後に一意の点列を持つ。
$p$ から出る右点加群を $M$ とする。$M$ と各終点単純の左Koszul候補複体とのテンソル積は、
長さ $n$ で右から $\min(n,4)+1$ 項だけを持つ。全項があれば点列上の階数は $1,3,3,1$ で完全である。
従って正の第一、第二ホモロジーはそれぞれ内部次数1,2のみに三次元ずつある。
極小解消の初めの二写像は候補と一致し、項3は全ての三次重なりを含むので、
真の第一Torはこの第一ホモロジーに等しく、第二Torは第二ホモロジーの商である。
よって $M=e_iA/H_iA$ であり、$H_i\otimes e_{i+1}A\to e_iA$ の関係は次数2に高々三個で生成される。
実際 $\bigwedge^2H$ の三つはそこに独立にあるので、ちょうどこれらが生成する。

抽象多項式代数 $P$ のKoszul複体から、直前の完全性は
$\operatorname{Tor}^P_1(k_i,A)=0$ を意味する。各終点の有界下次数付き加群の極小表示を取れば
左 $P$ 自由である。始点がどこでも商 $M$ は各長さで一次元なので、各終点で自由生成元は
各長さに一個あり、$\dim A_{i,i+n}=\binom{n+3}{3}$ を得る。次数零の持上げも自由生成元であり
$P$ 自身が埋まる。逆向きの点列で右自由性も従う。

右単純から $A$ へのHomは、非零な一文字による右 $P$ 乗法が単射なので零である。
正次数根基 $e_iA_+$ から $A$ への任意の斉次右加群写像を $e_iP_+$ に制限する。
$A$ が右 $P$ 自由で $P$ 上の単純の第一Extが $P$ に対して零なので、
その制限はある元の左乗法である。その左乗法を引いた写像は $H_i A$ を消し、
点加群の正次数尾部を経由する。この尾部から $A$ への加群写像は零である。
実際、各点加群の斉次元は次の $H$ の少なくとも二次元の部分で消えるが、
自由な右 $P$ 加群の非零な元を消す非零一次式はない。
よって第一Extも零である。反対側でも同じ二消滅を得る。候補Koszul複体の一端の完全性は
生成・二次表示、他端はFrobenius双対とこの二消滅から従い、中央はHilbert関数による。
\end{proof}

\begin{theorem}\label{thm:septic-point-graph}
定理\ref{thm:point-graph-curve-degree}の上限7に達する場合、点グラフ $X$ は純曲線で、
算術種数は5である。特に $X$ が滑らかな整曲線なら、全ての有限Frobenius候補が全次数で正則である。
その添字保持同型類は次の有限条件を満たす順序付きデータ $(X,p,q)$ の曲線同型類で尽きる。
\begin{enumerate}
\item $X$ は種数5の滑らかな射影曲線で $g^1_3$ を持たず、$p-q$ はちょうど位数4のPicard類を持つ。
\item $L=\omega_X(-p),\ N=\omega_X(-q)$ の完全線型系は閉埋込みで、
$R=\ker(H^0(L)\otimes H^0(N)\to H^0(LN))$ を取る。
$R,R^{\rm op}$ の交互窓に有限双対Frobenius検査を施し、全辺で射影ファイバーが空か一点、
一辺の像が一次元であることを課す。
\end{enumerate}
最小周期は $p,q$ を交換する曲線自己同型があれば1、それ以外は2である。
この記述は有限条件の解の存在やその次元を仮定しない。
\end{theorem}
\begin{proof}[残余三次曲線]
前節の五式の完全交叉 $C$ において、純曲線部分 $Y$ の残余 $Z$ は次数3で
$h^1(\mathcal O_Z)>0$、両側の射影三空間に埋まる。そのような純三次曲線は平面三次曲線である。
実際、一般の直線への線型射影 $f$ は有限平坦であり、跡で単位を分裂させると
$f_*\mathcal O_Z=\mathcal O\oplus\mathcal O(a)\oplus\mathcal O(b)$、$a\ge b$, $b\le-2$ となる。
有限平坦性は純性により直線上無捩れであることから従う。残りの二つの空間座標は、
射影直線の局所座標で割るとこの代数を生成する。それらの跡零成分は
$\mathcal O(-1)$ からの写像であり、$\mathcal O(b)$ 成分は零なので $a\ge-1$。
$\mathcal O(a)$ の平方が $\mathcal O(b)$ に非零成分を持たなければ二成分だけの部分代数に留まるため
$2a\le b$ である。従って $(a,b)=(-1,-2)$、残り二座標の跡零成分は比例する。
これが空間の一次式を与え、純性から平面内の三次因子である。

随伴により $\omega_Z=\mathcal O_Z$ かつ $h^0(\mathcal O_Z)=1$。
第六式は $I_{Y/C}(1,1)=\omega_Z$ の生成元である。従って六式は $Y$ をスキームとして切り出し、
$X=Y$ である。残余の両側次数は3なので
$\chi(I_{Y/C})=\chi(\mathcal O_Z(-1,-1))=-6$ であり、$\chi(\mathcal O_X)=-10+6=-4$ を得る。
\end{proof}

\begin{proof}[滑らかな場合の偏極と関係]
次数7の滑らかな種数5の空間曲線は平面にも二次曲面にも含まれない。
平面では平面種数公式に反する。滑らかな二次曲面上なら二次数 $(a,7-a)$ に対して
$(a-1)(6-a)=5$ に整数解がない。二次錐の解消 $\mathbb F_2$ 上の狭義変換は
元の滑らかな曲線に同型である（射は双有理かつ有限）。負切断 $s$ とファイバー $f$ で
その類は $as+7f$、随伴による種数は $-a^2+7a-6$ でやはり5にならない。
各 $L_i$ は独立な四切断を持ち、Riemann--Rochで次数1の $\omega_X/L_i$ が有効だから
一意な点 $p_i$ で $L_i=\omega_X(-p_i)$、$h^0(L_i)=4$ となる。
次数3以下の鉛筆があれば、その動く部分の $p_i$ を通る因子に点を加えて
$p_i+D$、$\deg D=2,\ h^0(p_i+D)\ge2$ とできる。
Riemann--Rochで $h^0(L_i(-D))\ge3$ となり埋込みに反する。

偏極の三階差が零なので $(1-T)^3(1+T)=1-2T+2T^3-T^4$ を使えば
$p_i+2p_{i+3}\sim 2p_{i+1}+p_{i+4}$ である。
三次の相異なる有効因子が線型同値なら鉛筆になるので、両辺は因子として等しい。
重複点の比較により $p_{i+3}=p_{i+1}$、すなわち二位置で戻る。
全て同じ点なら空間曲線に二次式がないため六関係が全て可換となり像が三次元になる。
従って相異なる二点 $p,q$ が交互に現れ、三階差から $4p\sim4q$ となる。
次数2以下の鉛筆もないので位数はちょうど4である。

$U=H^0(\omega_X)$ を五次元の標準座標空間とする。
標準像を消す二次式の空間の次元は3である。実際 $\dim\operatorname{Sym}^2U=15$、
$h^0(\omega_X^2)=12$ で下限3を得る。点 $p$ での微分は三次元の余法空間に入り、
微分の消える二次式は $\operatorname{Sym}^2 H^0(\omega_X(-p))$ に属するから
上の空間二次式の非存在で零である。相異なる二つの超平面
$V_p,V_q\subset U$ の積の対称二乗での像は、$p,q$ で消える13次元の空間全体となる。
共通の三文字 $H=V_p\cap V_q$ と両側の補文字で単項式を書けば、補文字の二つの平方だけが欠ける。
上の標準像の三式は全てここに入り、曲線上の積の階数は10である。
従って六関係は全乗法核であり、その中に $\bigwedge^2H$ が入る。
全窓は記載した反転列となり、命題\ref{prop:commuting-three-point-graph}が正則性を与える。

逆に表の入力が有限検査を通れば、最初の点対応は指定の曲線を含む。
前節の次数上限と本節の純性から点グラフがその曲線に等しい
（純な曲線への有限台のイデアルは零）。同型は点グラフと各埋込みの偏極を保ち、
二点を順序付きで復元する。逆にこの二点を対応させる曲線同型は全ての切断乗法核を保つ。
一移動は二点の順序を交換し、二移動は元の表示を返す。
\end{proof}
